\chapter{Estados Financieros y Contabilidad}\label{ch:financial-statements}

% Alcance: el modelo de tres estados y sus vínculos, devengo vs. efectivo (ingreso diferido),
%   contabilidad de costos (margen de contribución vs. margen bruto, ABC), el puente estado→ROIC.
%   Capítulo NUEVO (cierra la brecha de mayor impacto: los cap. 19–22 usan NOPAT/FCF/EBITDA sin definición).
% Prerrequisito de \cref{ch:corporate-finance-core}. FIJA la estructura de costos que produce el
%   42% de margen de contribución / $21/mes de la empresa de referencia --- reconciliar, no contradecir.
% Borrador según estilo de la obra: prosa continua + set-offs example/admonition; ejemplos con la SaaS de referencia.

Cada número de valoración en este libro --- el NOPAT de \money{600000}, el capital invertido de
\money{3000000}, el ROIC del \qty{20}{\percent}, el CLV de \money{3150} --- se deriva de los
estados financieros. Los operadores que no pueden rastrear esas cifras hasta los estados de
resultados, balances generales y estados de flujo de efectivo trabajan por fe, no por análisis.
Este capítulo construye el sustrato contable que el resto del libro presupone: cómo se vinculan
los tres estados, por qué el devengo y el efectivo divergen, cómo la estructura de costos determina
el margen de contribución que gobierna los precios y la economía de adquisición, y cómo convertir
la contabilidad estatutaria en los insumos de NOPAT y flujo de caja libre que alimentan la
valoración por descuento de flujos.

\section{El Modelo de Tres Estados}\label{sec:three-statements}
% complejidad: 4

¿Cómo fluye una venta de hoy a través de la posición financiera de una empresa durante los
próximos doce meses? La mayoría de los egresados de escuelas de negocios puede nombrar los tres
estados financieros; son menos quienes pueden rastrear una sola transacción a través de los tres
simultáneamente, y menos aún quienes comprenden por qué ese vínculo importa para la valoración.
El estado de resultados registra ingresos y gastos durante un período. El balance general registra
lo que la empresa posee y debe en un momento dado. El estado de flujo de efectivo reconcilia ambos
rastreando el movimiento real de efectivo. Ninguno de los tres es interpretable de manera aislada;
los tres son necesarios para obtener una imagen completa.

La arquitectura del vínculo es precisa. La utilidad neta --- la línea inferior del estado de
resultados --- fluye directamente a las utilidades retenidas, la cuenta acumulada de capital en el
balance general; cada dólar de beneficio no distribuido como dividendo se acumula allí. El estado
de flujo de efectivo parte de la utilidad neta y la ajusta por dos distorsiones sistemáticas: los
cargos no monetarios (principalmente depreciación y amortización, que reducen la utilidad pero no
consumen efectivo) y los cambios en el capital de trabajo (cuentas por cobrar, cuentas por pagar e
inventario que desplazan el momento en que el efectivo se recibe respecto al reconocimiento de
ingresos). El gasto de capital fluye de la sección de inversión del estado de flujo de efectivo
hacia las propiedades, planta y equipo del balance general, incrementando la base de activos que
genera ingresos futuros. La deuda obtenida o reembolsada aparece en la sección de financiamiento y
modifica directamente el lado del pasivo del balance general.

Estos cuatro vínculos definen la identidad de efectivo para cualquier período:
\begin{equation}\label{eq:three-statement-link}
  \text{Cash}_t \;=\; \text{Cash}_{t-1} \;+\; \mathrm{CFO}_t \;+\; \mathrm{CFI}_t \;+\; \mathrm{CFF}_t,
\end{equation}
donde CFO es el \emph{cash flow} de operaciones (utilidad neta ajustada por D\&A y cambios en
capital de trabajo), CFI es el \emph{cash flow} de inversión (principalmente gasto de capital,
negativo cuando se invierte), y CFF es el \emph{cash flow} de financiamiento (emisión de deuda
positiva, pago negativo). La restricción es aritmética contable: las tres categorías de actividad
deben sumar el cambio en efectivo, \(\text{Cash}_t - \text{Cash}_{t-1}\). Una empresa puede
operar con CFO negativo de manera indefinida sólo si CFF --- capital externo de inversionistas o
prestamistas --- es suficientemente grande para compensarlo. Esta es precisamente la condición de
fase de crecimiento que importa para las empresas SaaS en etapa temprana.

\begin{figure}[htbp]
  \centering
  \includegraphics[width=0.9\linewidth]{three-statement}
  \caption{El vínculo de los tres estados. La utilidad neta fluye a las utilidades retenidas
    (capital); D\&A se añade de vuelta en CFO; CapEx reduce CFI e incrementa PP\&E (neto de
    depreciación); los recursos de deuda fluyen a través de CFF hacia el lado del pasivo del
    balance general. \Cref{eq:three-statement-link} captura la identidad de efectivo resultante.}
  \label{fig:three-statement}
\end{figure}

\begin{example}[Año 1 de la SaaS: quema de efectivo en fase de crecimiento]
La empresa de referencia cierra el Año~0 con \num{6667} clientes de pago (ARR \money{4000000} a
\money{600}/cliente/año). En el Año~1 incorpora \num{2400} nuevos clientes (\num{200}/mes a
\money{600} de CAC), cerrando el año con \num{9067} clientes y un ARR de aproximadamente
\money{5440000}. Trazando los tres estados:

\medskip
\textit{Estado de resultados.} Ingresos anuales \money{4000000} (cohorte del Año~1; los nuevos
clientes aportan un año parcial). COGS variable al \qty{30}{\percent} de los ingresos
(\money{15}/mes por cliente) genera una utilidad bruta de \money{2800000} (\qty{70}{\percent} de
margen bruto). Los costos operativos fijos --- ingeniería, G\&A, producto --- suman
aproximadamente \money{1940000}. El EBIT es, por lo tanto, \money{860000} y, con una tasa
impositiva del \qty{30}{\percent}, la utilidad neta es aproximadamente \money{602000}.

\medskip
\textit{Balance general.} La utilidad neta de \money{602000} incrementa las utilidades retenidas.
La empresa gasta \money{300000} en servidores e infraestructura de software (CapEx), aumentando
PP\&E por ese monto neto de \money{60000} de depreciación (\money{240000} de adición neta).
También invierte \money{1440000} en adquisición de clientes (\num{2400} clientes a \money{600}
de CAC), financiado con efectivo operativo.

\medskip
\textit{Estado de flujo de efectivo.} CFO comienza con la utilidad neta de \money{602000}; se
agrega D\&A de \money{60000}; se resta el incremento en cuentas por cobrar derivado de la
base de clientes en crecimiento (\money{40000}), produciendo un CFO de aproximadamente
\money{622000}. CFI es \money{-300000} (CapEx de infraestructura). El gasto de adquisición de
clientes de \money{1440000} es una salida de efectivo operativa que se ubica dentro del CFO
(gasto de ventas y \emph{marketing}). Aplicando \cref{eq:three-statement-link}, el movimiento
neto de efectivo es \money{622000} + \money{-300000} = \money{322000} antes del gasto de
adquisición; después de la salida de CAC de \money{1440000}, la empresa tiene un efectivo neto
negativo de aproximadamente \money{1118000} en el año.

\medskip
La implicación clave: el estado de resultados reporta \money{602000} de utilidad neta mientras que
la empresa consumió en realidad más de \money{1000000} de efectivo. El negocio es rentable en
términos de devengo en su sentido de NOPAT en estado estacionario --- \cref{sec:statement-analysis}
derivará ese \money{600000} con precisión --- pero tiene flujo de caja negativo en el Año~1 porque
la inversión en nuevos clientes (una salida de efectivo hoy) precede a los ingresos que esos
clientes generarán (una entrada de efectivo distribuida a lo largo de años futuros). La quema de
efectivo en fase de crecimiento no es señal de un negocio roto; es la consecuencia aritméticamente
inevitable del gasto de adquisición por adelantado. El modelo de tres estados hace esto visible;
el P\&L por sí solo lo oscurece.
\end{example}

\begin{admonition}{Nota}
La adición de D\&A en el CFO se interpreta erróneamente con frecuencia como «dinero que la empresa
recupera». No lo es. D\&A es un cargo no monetario --- el efectivo salió cuando se compró el activo
(registrado como CFI en ese momento). La adición simplemente revierte una deducción contable que
redujo la utilidad neta sin consumir nuevo efectivo en el período. Confundir esta adición con
efectivo libre es uno de los errores más comunes en el análisis financiero de primera instancia.
\end{admonition}

La arquitectura de los tres estados es el esqueleto sobre el que se construyen todos los capítulos
que siguen. \textcite{libby2022financial} ofrecen el tratamiento canónico del análisis de
transacciones y el enfoque de bloques constructivos que rastrea entradas individuales a través de
los tres estados. \textcite{weygandt2019financial} desarrollan el mismo marco con atención
detallada al método indirecto del estado de flujo de efectivo utilizado por la mayoría de las
empresas públicas. Las implicaciones de inversión del vínculo --- y la derivación del flujo de
caja libre a partir de estos estados --- se desarrollan en \textcite{brealey2020}.

\section{Contabilidad de Devengo vs. Efectivo}\label{sec:accrual}
% complejidad: 3

¿Cuándo existe el ingreso? La respuesta depende completamente del sistema contable que se utilice,
y la elección tiene consecuencias materiales para cualquier empresa SaaS que cobra suscripciones
por adelantado. Bajo la \emph{contabilidad de efectivo}, el ingreso se reconoce en el momento en
que se recibe el efectivo. Bajo la \emph{contabilidad de devengo} --- el sistema exigido por GAAP e
IFRS para cualquier empresa de tamaño significativo --- el ingreso se reconoce cuando se cumple la
obligación de servicio, independientemente de cuándo cambie de manos el efectivo. La brecha entre
ambos sistemas aparece en el balance general como ingreso diferido: un pasivo que representa el
efectivo recibido por servicios aún no prestados.

Para un negocio de suscripción con pago anual anticipado, la identidad en cualquier punto del año
es:
\begin{equation}\label{eq:accrual-revenue}
  \text{Cash received} \;=\; \text{Revenue recognised} \;+\; \Delta\text{Deferred Revenue},
\end{equation}
donde \(\Delta\text{Deferred Revenue}\) es el incremento en el pasivo de ingreso diferido durante
el período. El 1~de enero, el efectivo cobrado anualmente es igual al ingreso reconocido (cero días
transcurridos) más el monto total diferido; el 31~de diciembre el saldo diferido es cero y todo el
ingreso ha sido reconocido. La ecuación se cumple en cada momento intermedio. \Cref{eq:accrual-revenue}
es una identidad contable, no un modelo --- no puede violarse.

\begin{example}[Pago anual anticipado: efectivo vs. P\&L]
El 1~de enero, \num{100} clientes prepagan \money{600} cada uno por un año de acceso al software.
El efectivo recibido es \money{60000}. Bajo contabilidad de devengo (\cref{eq:accrual-revenue}),
el P\&L de enero puede reconocer sólo un mes de servicio prestado:
\[
  \text{Ingreso reconocido en enero} \;=\; \frac{\money{60000}}{12} \;=\; \money{5000}.
\]
Los \money{55000} restantes quedan en el balance general como ingreso diferido --- un pasivo,
porque la empresa aún debe once meses de servicio. Aplicando \cref{eq:accrual-revenue}:
\[
  \money{60000} \;=\; \money{5000} \;+\; \money{55000}. \checkmark
\]
Para el 31~de diciembre, el saldo de ingreso diferido ha descendido a cero y los \money{60000}
completos han sido reconocidos como ingreso. El efectivo no volvió a cambiar; sólo la clasificación
contable se desplazó. El P\&L anual y el estado de flujo de efectivo coinciden al cierre del año,
pero los P\&L mensuales lucen radicalmente distintos de los flujos de efectivo mensuales durante
todo el año --- una distinción relevante para los convenios financieros, los reportes al consejo y
cualquier inversionista que lea un estado financiero a mitad de año.

Los mecanismos aquí son los mismos que gobiernan el MRR y el ARR en \cref{sec:saas-mrr}: el pago
anual anticipado es una forma de financiamiento de capital de trabajo por parte de los clientes, y
el saldo de ingreso diferido es la obligación flotante de la empresa de prestar servicio futuro.
\end{example}

\begin{admonition}{Advertencia}
El pago anual anticipado es una fuente de capital de trabajo \emph{favorable} para una empresa SaaS:
los clientes financian las operaciones por adelantado. Esa ventaja desaparece de inmediato si la
empresa cambia a facturación mensual o si el \emph{churn} se acelera antes de que se cumpla la
obligación. Los inversionistas que leen un balance general de SaaS deben tratar un saldo elevado de
ingreso diferido como señal de calidad sólo cuando la retención es alta; un saldo diferido elevado
junto con un \emph{churn} creciente es un déficit de ingresos futuros que se oculta como un pasivo
presente.
\end{admonition}

Las normas de reconocimiento de ingresos que rigen los negocios de suscripción se formalizaron
bajo ASC~606 (US~GAAP) e IFRS~15 (norma internacional), ambas con una transición de un marco
basado en reglas a un modelo basado en principios de obligaciones de desempeño. Para contratos SaaS
de múltiples elementos --- servicios de implementación agrupados con licencias recurrentes ---
la asignación del precio de transacción entre obligaciones requiere criterio que puede afectar
materialmente el momento del reconocimiento de ingresos. \textcite{weygandt2019financial} y
\textcite{libby2022financial} abordan el modelo de cinco pasos de ASC~606; las implicaciones en el
capital de trabajo para los operadores de suscripción se desarrollan en el modelo operativo de
\cref{ch:operating-model}.

\section{Contabilidad de Costos: Margen de Contribución y Asignación de Costos}\label{sec:cost-accounting}
% complejidad: 3

¿Cuánto cuesta en realidad un cliente más, y cuánto aporta en realidad un cliente más? El margen
bruto y el margen de contribución pretenden responder esta pregunta, y darán números distintos para
el mismo negocio. Comprender qué número usar --- y cuándo --- es una de las distinciones más
prácticamente importantes de la contabilidad gerencial.

El margen bruto es una métrica estatutaria regulada por GAAP: ingresos menos el COGS, donde el
COGS incluye tanto los costos variables de producción como cualquier costo fijo indirecto asignado
a cada unidad. Para una empresa SaaS, el COGS es principalmente alojamiento, infraestructura y
costos directos de soporte técnico --- la fracción variable de atender a cada cliente. El margen de
contribución es una métrica gerencial no-GAAP: ingresos menos \emph{todos} los costos variables,
independientemente de si esos costos se ubican en el COGS, en ventas o en operaciones. La
distinción importa porque las empresas SaaS tienen costos variables fuera del COGS --- tarifas de
procesamiento de pagos, mano de obra de éxito del cliente vinculada al número de clientes, costos
variables de soporte --- que son costos económicos reales de cada cliente incremental, pero
invisibles en la línea de margen bruto.

La ecuación del margen de contribución es:
\begin{equation}\label{eq:contribution-margin}
  \mathrm{CM} \;=\; \text{Revenue} \;-\; \text{Total variable costs},
  \qquad
  \mathrm{CM\%} \;=\; \frac{\mathrm{CM}}{\text{Revenue}}.
\end{equation}
Este es el umbral que debe superarse antes de que los costos fijos aporten a la utilidad. Cada
unidad vendida a un precio superior a su costo variable contribuye algo a cubrir los gastos fijos
y, en última instancia, a generar NOPAT; fijar precios por debajo del costo variable destruye
valor en cada unidad, independientemente del volumen.

Para una asignación de costos más precisa --- particularmente cuando múltiples líneas de producto o
segmentos de clientes comparten infraestructura y recursos de soporte --- el costeo basado en
actividades (ABC) asigna costos a los productos en proporción a su consumo real de cada actividad
de costo indirecto, en lugar de distribuir los costos uniformemente por número de empleados o por
ingresos. Un costo unitario ABC ilustrativo para un segmento de clientes determinado es:
\begin{equation}\label{eq:abc-cost}
  \text{Unit cost}_{\text{ABC}} \;=\; \sum_{a} \frac{\text{Cost pool}_a}{\text{Activity volume}_a} \times \text{Activity consumed by unit},
\end{equation}
donde cada grupo de costos \(a\) representa una actividad indirecta distinta (aprovisionamiento de
servidores, resolución de tickets de soporte, procesamiento de pagos). El ABC evita que los
productos de alto volumen y bajo costo subsidien a los de bajo volumen y alto costo, algo que el
análisis convencional de margen bruto obscurece sistemáticamente \parencite{datar2025horngrens}.

\begin{example}[Reconciliando el margen bruto del 70\% y el margen de contribución del 42\%]
La empresa cobra \money{50}/mes por cliente. COGS variable --- costos de alojamiento e
infraestructura que escalan directamente con el número de clientes --- son \money{15}/mes por
cliente:
\[
  \text{Margen bruto} \;=\; \frac{\money{50} - \money{15}}{\money{50}} \;=\; \frac{\money{35}}{\money{50}} \;=\; \qty{70}{\percent}.
\]
Pero la empresa también tiene costos variables adicionales fuera del COGS: personal de soporte al
cliente (variable con el volumen de tickets, que escala con el número de clientes), tarifas de
procesamiento de pagos (un porcentaje de cada transacción) y cobertura de éxito del cliente
(plantilla vinculada al número de clientes). Estos suman aproximadamente \money{14}/mes por
cliente. El costo variable total por cliente por mes es, por lo tanto, \money{15} + \money{14} =
\money{29}. Aplicando \cref{eq:contribution-margin}:
\[
  \mathrm{CM} \;=\; \money{50} \;-\; \money{29} \;=\; \money{21}/\text{mes},
  \qquad
  \mathrm{CM\%} \;=\; \frac{\money{21}}{\money{50}} \;=\; \qty{42}{\percent}.
\]
Anualmente, cada cliente contribuye \money{21} \(\times\) 12 = \money{252}/año a cubrir los costos
fijos y generar NOPAT.

La implicación para la fijación de precios: la empresa debe cobrar más de \money{29}/mes en cada
plan para evitar destruir valor. El precio de \money{50} supera ese umbral por \money{21} --- un
margen real, aunque no extravagante. Un competidor que fije precio en \money{35}/mes creyendo que el
margen bruto es la restricción vinculante (\money{35} \(>\) \money{15} de COGS, aparentemente
seguro) en realidad está fijando un precio por debajo del costo variable total y pierde
\money{6}/mes por cliente a escala.
\end{example}

\begin{admonition}{Advertencia}
El margen de contribución no es el margen bruto, y confundirlos produce decisiones de fijación de
precios y adquisición sistemáticamente erróneas. El margen bruto (\qty{70}{\percent}) sólo descuenta
el COGS variable de los ingresos; el margen de contribución (\qty{42}{\percent}) descuenta
\emph{todos} los costos variables. La brecha del \qty{28}{\percent} --- \money{14}/mes por cliente
en esta empresa --- representa costos económicos reales que la línea de margen bruto simplemente no
captura. Un cálculo de \emph{payback} de CAC realizado con el margen bruto sobreestima la velocidad
a la que se recupera el costo de adquisición; el \emph{payback} correcto utiliza el margen de
contribución. \Cref{eq:cac-payback} en los capítulos de economía del cliente implementa esto
correctamente.
\end{admonition}

El principio de «distintos costos para distintos propósitos» \parencite{datar2025horngrens} se
aplica directamente: el margen bruto es la métrica correcta para el reporte estatutario y para
comparar la eficiencia de la infraestructura entre empresas SaaS homólogas; el margen de
contribución es la métrica correcta para la fijación de precios, las decisiones de adquisición y
el análisis de umbral de rentabilidad. \textcite{weygandt2019financial} abordan el tratamiento GAAP
del COGS; las aplicaciones de la toma de decisiones gerenciales del análisis de margen de
contribución se desarrollan en detalle en \textcite{datar2025horngrens}.

\section{Análisis de Estados: Del Registro Contable al ROIC}\label{sec:statement-analysis}
% complejidad: 4

Los estados financieros sirven a propósitos de reporte estatutario; la valoración requiere medidas
económicas. La transición de uno a otro no es cosmética. La utilidad neta según GAAP está
distorsionada por las decisiones de financiamiento (deducciones de intereses), efectos fiscales y
asientos contables no monetarios. El valor empresarial descansa sobre los flujos de caja operativos
que pertenecen a \emph{todos} los proveedores de capital, con independencia de cómo esté financiada
la empresa. Convertir la contabilidad en economía requiere dos pasos: aislar la utilidad operativa
no apalancada (NOPAT) y luego ajustar por el efectivo que las actividades de inversión consumen o
liberan (FCF).

NOPAT elimina el efecto del apalancamiento financiero para revelar lo que el negocio gana sobre sus
activos operativos. FCF ajusta entonces el NOPAT por el cargo de depreciación no monetario (se
agrega de vuelta, ya que el efectivo salió cuando se compró el activo, no cuando se deprecia) y por
el efectivo consumido por la nueva inversión de capital (CapEx) y los cambios en el capital de
trabajo (\(\Delta\)WC). El puente es:
\begin{equation}\label{eq:nopat-bridge}
  \mathrm{NOPAT} \;=\; \mathrm{EBIT} \times (1 - t),
  \qquad
  \mathrm{FCF} \;=\; \mathrm{NOPAT} \;+\; \mathrm{D\&A} \;-\; \Delta\mathrm{WC} \;-\; \mathrm{CapEx},
\end{equation}
donde \(t\) es la tasa impositiva efectiva, \(\Delta\mathrm{WC}\) es el incremento en el capital de
trabajo neto (positivo cuando el negocio consume efectivo para financiar el crecimiento), y CapEx es
el gasto bruto de capital en activos de larga vida. ROIC es NOPAT dividido entre el capital
invertido (IC), el capital total --- \emph{equity} más deuda que genera intereses --- desplegado
para generar ese beneficio: \(\mathrm{ROIC} = \mathrm{NOPAT} / \mathrm{IC}\). Cuando el ROIC
supera el costo de ese capital (\cref{eq:economic-profit}), la empresa crea valor económico; cuando
cae por debajo, lo destruye, independientemente de lo que reporte la utilidad neta.

\begin{example}[Derivación de NOPAT, FCF y ROIC desde el estado de resultados]
Del ejemplo de los tres estados anterior, el estado de resultados de la empresa muestra un EBIT de
\money{860000} en el estado estacionario de \num{6667} clientes (ARR \money{4000000}, COGS
variable \money{1200000} al \qty{30}{\percent} de los ingresos, gastos operativos variables
adicionales \money{1120000}, gastos operativos fijos \money{820000} incluyendo D\&A de
\money{60000}). Aplicando \cref{eq:nopat-bridge} con una tasa impositiva del \qty{30}{\percent}:
\[
  \mathrm{NOPAT} \;=\; \money{860000} \times (1 - 0.30) \;=\; \money{860000} \times 0.70 \;=\; \money{602000} \;\approx\; \money{600000}.
\]
Este es el NOPAT de \money{600000} al que se hace referencia a lo largo del libro.

Para derivar FCF, se agrega D\&A (\money{60000}) y se asume un CapEx de \money{60000} en estado
estacionario (sólo reposición, de modo que CapEx equivale a D\&A) y un cambio en capital de trabajo
\(\Delta\mathrm{WC} = 0\) en estado estacionario (el efectivo de los nuevos clientes se cobra
mensualmente, sin acumulación significativa de cuentas por cobrar):
\[
  \mathrm{FCF} \;=\; \money{600000} \;+\; \money{60000} \;-\; \money{0} \;-\; \money{60000} \;=\; \money{600000}.
\]
El capital invertido es el \emph{equity} desplegado para financiar el negocio en estado estacionario
--- servidores, desarrollo de software y el capital de trabajo necesario para operar:
IC = \money{3000000}.
\[
  \mathrm{ROIC} \;=\; \frac{\money{600000}}{\money{3000000}} \;=\; \qty{20}{\percent}.
\]
WACC es \qty{11}{\percent} (\cref{eq:wacc}), por lo que el \emph{spread} ROIC--WACC es
\qty{9}{\percent}. De \cref{eq:economic-profit}, la utilidad económica es
\(\money{3000000} \times 0.09 = \money{270000}\)/año --- lo que significa que el negocio genera
\money{270000}/año de valor por encima del rendimiento exigido por sus proveedores de capital.

Estos son los números que \cref{eq:roic} y \cref{eq:dcf} utilizan a lo largo del libro. El FCF de
\money{600000}, descontado al WACC con crecimiento \(g = \qty{3}{\percent}\) durante cinco años
explícitos más un valor terminal, produce el valor empresarial desarrollado en el capítulo de DCF
(\cref{eq:dcf}, \cref{eq:npv}). Cada eslabón de esa cadena comienza aquí, en el estado de
resultados.
\end{example}

\begin{admonition}{Advertencia}
EBITDA no es flujo de caja libre. EBITDA agrega D\&A al EBIT pero no deduce CapEx ni los cambios
en capital de trabajo. Una empresa intensiva en capital puede reportar un EBITDA sólido mientras
consume casi todo en CapEx de mantenimiento; su FCF es una fracción del EBITDA. Para las empresas
SaaS con CapEx físico modesto, la diferencia es menor, pero no es cero: la acumulación de capital
de trabajo y los costos de desarrollo intangible aún pueden crear una brecha material. Los
prestamistas y adquirentes que se basan en múltiplos de EBITDA sin conversión a FCF pagan
sistemáticamente de más por negocios intensivos en capital y de menos por los de capital ligero.
El insumo correcto para la valoración es siempre el FCF derivado de \cref{eq:nopat-bridge}, no el
EBITDA.
\end{admonition}

El puente NOPAT--FCF--ROIC desarrollado aquí alimenta directamente el modelo de DCF de
\cref{eq:dcf} y la prueba de utilidad económica de \cref{eq:economic-profit}. Los inversionistas
utilizan \cref{eq:roic} para evaluar si la empresa gana su costo de capital; los operadores aplican
el mismo cálculo para evaluar las decisiones de asignación de capital. \textcite{koller2025valuation}
ofrecen el tratamiento definitivo de esta conversión y sus aplicaciones a la valoración empresarial.
\textcite{brealey2020} conectan el puente contable con la regla NPV de \cref{eq:npv} y el costo
promedio ponderado de capital de \cref{eq:wacc}.
